\begin{textsample}{POS Dim 3 – ac – Score 48.00 – ac/ac\_inf\_25.txt}  \label{ex:f3_pos_028}
25. \textbf{CRIANÇAS} \textbf{PEQUENAS} ( 4 ANOS A 5 ANOS E 11 \textbf{MESES} ) 25.1.
Campo de experiências “ \textbf{Escuta}, fala, pensamento e imaginação ” Desde o \textbf{nascimento}, as \textbf{crianças} participam de situações comunicativas cotidianas com as pessoas com as quais interagem.
As primeiras formas de interação do \textbf{bebê} são os movimentos do seu corpo, o olhar, a postura corporal, o sorriso, o choro e outros recursos vocais, que ganham sentido com a interpretação do outro.
Progressivamente, as \textbf{crianças} vão ampliando e enriquecendo seu vocabulário e demais recursos de expressão e de compreensão, apropriando -se da língua materna – que se torna, pouco a pouco, seu veículo privilegiado de interação.
Na Educação \textbf{Infantil}, é importante promover experiências nas quais as \textbf{crianças} possam falar e \textbf{ouvir}, potencializando sua participação na cultura oral, pois é na \textbf{escuta} de histórias, na participação em conversas, nas descrições, nas narrativas elaboradas individualmente ou em grupo e nas implicações com as múltiplas linguagens que a \textbf{criança} se constitui ativamente como sujeito singular e pertencente a um grupo social.
Desde cedo, a \textbf{criança} manifesta \textbf{curiosidade} com relação à cultura escrita : ao \textbf{ouvir} e acompanhar a leitura de textos, ao observar os muitos textos que circulam no contexto familiar, comunitário e escolar, ela vai construindo sua concepção de língua escrita, reconhecendo diferentes usos sociais da escrita, dos gêneros, suportes e portadores.
Na Educação \textbf{Infantil}, a \textbf{imersão} na cultura escrita deve partir do que as \textbf{crianças} conhecem e das \textbf{curiosidades} que deixam transparecer.
As experiências com a literatura \textbf{infantil}, propostas pelo educador, mediador entre os textos e as \textbf{crianças}, contribuem para o desenvolvimento do gosto pela leitura, do \textbf{estímulo} à imaginação e da ampliação do conhecimento de mundo.
Além disso, o contato com histórias, contos, fábulas, poemas, cordéis etc. propicia a familiaridade com \textbf{livros}, com diferentes gêneros literários, a diferenciação entre ilustrações e escrita, a aprendizagem da direção da escrita e as formas corretas de manipulação de \textbf{livros}.
Nesse convívio com textos escritos, as \textbf{crianças} vão construindo hipóteses sobre a escrita que se revelam, inicialmente, em rabiscos e garatujas e, à medida que vão conhecendo letras, em escritas espontâneas, não convencionais, mas já \textbf{indicativas} da compreensão da escrita como sistema de representação da língua. 25.2.
Direitos de aprendizagem e desenvolvimento - Conviver com \textbf{crianças} e \textbf{adultos} em situações comunicativas do cotidiano, desenvolvendo progressivamente confiança em participar de experiências que envolvam o pensamento, a imaginação, sentimentos, narrativas e diálogos, quer seja em momentos individuais ou coletivos. - \textbf{Brincar} e se deliciar com histórias e poesias, com jogos, parlendas, rimas, \textbf{brinquedos} cantados, textos de imagens e escritos, ampliando seu repertório das manifestações culturais, favorecendo sua progressiva aprendizagem dos vários gêneros e formas de expressão : gestual, oral, escrita, \textbf{plástica}, \textbf{dramática} e musical. - Explorar \textbf{gestos}, expressões, rimas, imagens, canções, enredos de histórias, textos que compõem o patrimônio cultural da humanidade ( histórias, lendas, mitos, fábulas, poemas, parlendas, trava-línguas, piadas, adivinhas, músicas etc ), atuando como autores no processo, redescobrindo e transformando à sua maneira e de acordo com suas possibilidades, considerando o momento de seu desenvolvimento. - Participar de \textbf{rodas} de conversa, relatos de experiência, leitura e contação de histórias, saraus, dramatizações, apresentações, construção de narrativas, representação de papeis no faz de conta, exploração de materiais impressos e de variedade linguística, comunicando -se com os outros, trocando informações e agindo no mundo, criando assim a possibilidade de compreensão, de organização de ideias e consciência de si mesmo. - Expressar sentimentos, ideias, \textbf{desejos}, necessidades, dúvidas, informações e descobertas, possibilitando o compartilhamento de saberes e conhecimentos entre todos os envolvidos, sendo sua autoria garantida pelo \textbf{gesto}, pela fala, pelo desenho, pela escrita ( convencional ou não convencional ) ou por outra forma de \textbf{registro}. - Conhecer -se como sujeitos de sua cultura, reconhecendo e dando sentido às intenções comunicativas e de \textbf{escuta}, aos \textbf{gestos}, as emissões sonoras, as narrativas de diferentes gêneros linguísticos, aos diferentes suportes e gêneros textuais, valorizando suas origens. 25.2.1.
Objetivo de aprendizagem e desenvolvimento : ( EI03EF01 ) - Expressar ideias, \textbf{desejos} e sentimentos sobre suas vivências, por meio da linguagem oral e escrita ( escrita espontânea ), de \textbf{fotos}, desenhos e outras formas de expressão.
Aprendizagens específicas - Comunicar -se por meio da fala, ampliando progressivamente, sua capacidade de \textbf{ouvir} e de adequar a linguagem oral às diversas situações comunicativas do cotidiano. - Desenvolver a capacidade comunicativa, formulando perguntas, pedindo explicações, relatando acontecimentos, expressando suas ideias, sentimentos, dúvidas, opiniões em situações interativas. - Desenvolver a capacidade de \textbf{escuta} \textbf{atenta}, respeitando o ponto de vista do outro. - Desenvolver estratégias de comunicação para compreender e se fazer compreender, por meio de diferentes linguagens ( corporal, musical, visual, \textbf{plástica}, oral e escrita ), ajustadas às diferentes intenções e situações de comunicação. - Desenvolver o interesse e o gosto por comunicar suas ideias e opiniões, fazendo uso da escrita espontânea. - Apropriar -se, gradativamente, dos diversos usos da linguagem oral para \textbf{conversar}, \textbf{brincar}, comunicar suas experiências, expressar \textbf{desejos}, necessidades, opiniões, ideias, preferências e sentimentos, nas diversas situações de interações presentes no cotidiano. \textbf{Sugestões} de experiências - \textbf{Conversar} sobre variados temas. - Prever o final de uma história ou o que acontecerá com um determinado \textbf{personagem}. - Relatar experiências próprias ou de outras pessoas. - Compartilhar opiniões, argumentar e buscar soluções sobre diversas situações. - \textbf{Escutar} experiências e opiniões dos \textbf{colegas}. - Expressar -se por meio das diversas linguagens. - \textbf{Contar} como foram realizadas produções individuais ou coletivas. - Falar sobre as etapas de uma tarefa concluída ou a ser realizada. - Explicar e argumentar suas ideais, construindo situações de diálogo. 25.2.2.
Objetivo de aprendizagem e desenvolvimento : ( EI03EF02 ) - \textbf{Inventar} \textbf{brincadeiras} cantadas, poemas e canções, criando rimas, aliterações e ritmos.
Aprendizagens específicas - Perceber e explorar a sonoridade das palavras, nas \textbf{brincadeiras} e jogos orais ( adivinhas, trava-línguas, parlendas, brincos, quadrinhas e outras mnemônias ). - Ampliar a oralidade, enriquecendo o seu vocabulário por meio de músicas, narrativas ( poemas, histórias, contos, parlendas, conversas ), \textbf{brincadeiras} e jogos. - Desenvolver, progressivamente, o \textbf{hábito} e o prazer por \textbf{escutar}, recitar e ler textos poéticos, rimas, aliterações e ritmos. - Desenvolver a imaginação e a criatividade na produção de rimas, ritmos e aliterações, em \textbf{brincadeiras} e jogos. - Construir noções da linguagem oral e escrita, percebendo as relações de semelhanças e diferenças entre ambas. \textbf{Sugestões} de experiências - Recitar poemas, cantar músicas, ler textos diversos de forma não convencional. - \textbf{Brincar} com jogos orais envolvendo \textbf{sons}, ritmos e rimas. - Organizar e vivenciar saraus de poemas e músicas, com a participação de \textbf{crianças} de outras salas e famílias. - \textbf{Ouvir} recitação de poemas e jogos orais : adivinhas, trava-línguas, parlendas, brincos, quadrinhas e outros. - Apreciar espetáculos e apresentações musicais na \textbf{instituição} educativa e na comunidade. - Explorar \textbf{sons} semelhantes ou diferentes em atividades diversas. - \textbf{Brincar} com a sonoridade das palavras, explorando -as, refletindo e estabelecendo relações com suas escritas. - Explorar repertório de poemas e jogos orais conhecidos, de forma \textbf{lúdica} e prazerosa. 25.2.3.
Objetivo de aprendizagem e desenvolvimento : ( EI03EF03 ) - Escolher e folhear \textbf{livros}, procurando orientar -se por temas e ilustrações e tentando identificar palavras conhecidas.
Aprendizagens específicas - Desenvolver o \textbf{hábito}, o prazer e encanto pela leitura por meio da participação em diversas situações de \textbf{escuta} de histórias, recontos e narrativas realizadas por \textbf{adultos} e \textbf{crianças}. - Desenvolver o interesse pelo \textbf{manuseio} de \textbf{livros} de diversos gêneros textuais, familiarizando -se com diferentes autores, ilustradores e suportes. - Aprimorar sua sensibilidade estética para ler, mesmo que de forma não convencional, em momentos individuais ou coletivos. - Desenvolver a \textbf{curiosidade} e interesse pela pesquisa e busca de informações em \textbf{livros} diversos, \textbf{apoiando} -se na memória de textos e histórias conhecidas, em ilustrações e imagens. - Desenvolver, progressivamente, o comportamento leitor, por meio do \textbf{manuseio} de \textbf{livros}, revistas e outros portadores textuais. - Desenvolver a \textbf{percepção} das diferentes estruturas textuais, temas, ilustração e diagramação, explorando diversos gêneros literários que circulam socialmente. - Conhecer \textbf{livros} de autores e ilustradores locais, que retratam a vida na floresta, mitos e lendas regionais, indígenas, dentre outras. \textbf{Sugestões} de experiências - Folhear \textbf{livros}, revistas, histórias em quadrinhos e outros materiais escritos. - \textbf{Manusear} \textbf{livros} e materiais impressos \textbf{demonstrando} \textbf{cuidados}. - Escolher histórias e \textbf{livros} de seu interesse para \textbf{manusear} e ler à sua maneira. - Visitar espaços literários. - \textbf{Ouvir} e apreciar gêneros literários da cultura local e regional. - Pesquisar informações de temas variados em \textbf{livros}, revistas e outros materiais impressos. - Ler, mesmo que de forma não convencional, textos diversos, \textbf{apoiando} -se em imagens e ilustrações. - Explorar materiais impressos, observando suas diferentes diagramações. - Vivenciar situações de leitura com diversos tipos de textos, diferentes estruturas e tramas. 25.2.4.
Objetivo de aprendizagem e desenvolvimento : ( EI03EF04 ) - Recontar histórias \textbf{ouvidas} e planejar coletivamente roteiros de vídeos e de encenações, definindo os contextos, os \textbf{personagens}, a estrutura da história.
Aprendizagens específicas - Recontar histórias e contos conhecidos, recuperando características do enredo original, como descrição de \textbf{personagens}, cenários e objetos, com ou sem a \textbf{ajuda} do professor. - Ampliar o seu repertório de histórias preferidas, utilizando -o como suporte para a produção de recontos e construção de roteiros de vídeos ou encenações. - Conhecer a estrutura narrativa de histórias ( nacional, regional e local ) identificando \textbf{personagens}, cenários, tema, sequência cronológica, ação e intenção dos \textbf{personagens}. - Desenvolver a criatividade e a imaginação por meio da produção de reconto e planejamento de roteiro de vídeos e encenações. \textbf{Sugestões} de experiências - \textbf{Dramatizar} histórias preferidas na sua e em outras turmas. - Narrar histórias e gravar em áudio ou vídeo. - \textbf{Conversar} sobre \textbf{livros} e histórias lidas. - Recontar histórias conhecidas, usando diferentes entonações para criar efeitos de suspense, imitação da voz dos \textbf{personagens} e outros. - Recontar histórias conhecidas, aproximando o relato da história original. - Recontar histórias \textbf{manuseando} recursos tecnológicos e midiáticos. - Planejar, com a \textbf{ajuda} do professor, roteiros de histórias para encenação. - \textbf{Ouvir} e apreciar histórias, identificando os elementos da narrativa, com \textbf{ajuda} do professor e dos \textbf{colegas}. 25.2.5.
Objetivo de aprendizagem e desenvolvimento : ( EI03EF05 ) - Recontar histórias \textbf{ouvidas} para produção de reconto escrito, tendo o professor como escriba.
Aprendizagens específicas - Desenvolver o interesse por \textbf{ouvir} a leitura de diferentes tipos de textos e de gêneros textuais, feita pelo ( a ) professor ( a ), desenvolvendo comportamento leitor. - Recontar histórias e contos conhecidos, preservando as características da linguagem do texto recontado. - Compreender, gradualmente, as relações entre a linguagem oral e linguagem escrita, percebendo semelhanças e diferenças. - Produzir texto coletivamente, tendo o professor como escriba, percebendo que as ideias se organizam em formas convencionais de escrita. - Produzir textos orais observando e mantendo a sequência dos fatos da narrativa. \textbf{Sugestões} de experiências - \textbf{Conversar} sobre \textbf{livros} lidos pelo professor ( texto e ilustrações ), sobre as suas produções e as dos \textbf{colegas}. - Recontar histórias e contos – individual e/ou coletivamente, tendo o professor como escriba. - Organizar o mural da turma com textos poéticos, notícias, elaborados coletivamente. - Ditar textos diversos, individual ou coletivamente, para o professor escrever. - \textbf{Conversar} sobre as características dos \textbf{personagens} e das histórias lidas e recontadas. 25.2.6.
Objetivo de aprendizagem e desenvolvimento : ( EI03EF06 ) - Produzir suas próprias histórias orais e escritas ( escrita espontânea ), em situações com função social significativa.
Aprendizagens específicas - Produzir, ainda que não convencionalmente, escritas em situações que se aproximem do uso social, de algumas histórias que conheça de memória e outras de autoria. - Desenvolver o interesse pela leitura de histórias da literatura \textbf{infantil} ( internacional, nacional, regional e local ) e de diferentes textos que circulam socialmente. - Apreciar a estrutura da escrita de algumas histórias, fábulas e contos. - Ampliar, gradativamente, o vocabulário, utilizando -o em suas narrativas e \textbf{brincadeiras}. - Desenvolver comportamento leitor e produtor de textos, levando em consideração situações que se aproximem do uso social. - Ampliar a oralidade por meio da construção de narrativas de suas próprias vivências. - Apropriar -se, gradualmente, dos aspectos gráficos da escrita, percebendo -a como elemento de comunicação e interação social. \textbf{Sugestão} de experiências - Criar histórias orais livremente. - Recontar histórias em diversas situações. - Narrar situações vivenciadas no cotidiano. - Escrever individual ou coletivamente textos diversos. - Escrever espontaneamente, utilizando diferentes ferramentas e suportes de escrita. 25.2.7.
Objetivo de aprendizagem e desenvolvimento : ( EI03EF07 ) - Levantar hipóteses sobre gêneros textuais veiculados em portadores conhecidos, recorrendo a estratégias de observação gráfica e/ou de leitura.
Aprendizagens específicas - Reconhecer diversos suportes e gêneros textuais da linguagem escrita, identificando alguns aspectos de sua estrutura gráfica. - Identificar alguns elementos nos \textbf{livros} de literatura \textbf{infantil} : capa, ilustração, título, autor, ilustrador, editora, \textbf{personagens}, dentre outros elementos. - Perceber os aspectos gráficos dos diversos gêneros textuais, em situações significativas de leitura e escrita. - Desenvolver postura crítica, refletindo sobre a escrita e a leitura, usando estratégias de observação gráfica e/ou de leitura. \textbf{Sugestões} de experiências - \textbf{Ouvir} a leitura de textos de diferentes gêneros. - Participar de situações de leitura, mesmo sem saber ler convencionalmente. - Identificar seu próprio nome e diferentes nomes de pessoas ou objetos. - Participar de \textbf{brincadeiras} e jogos em que reconheçam o próprio nome e dos nomes dos \textbf{colegas}. - Explorar com os \textbf{colegas} materiais impressos variados, de diferentes gêneros. - Identificar palavras em textos que sabe de memória. 25.2.8.
Objetivo de aprendizagem e desenvolvimento : ( EI03EF08 ) - Selecionar \textbf{livros} e textos de gêneros conhecidos para a leitura de um \textbf{adulto} e/ou para sua própria leitura ( partindo de seu repertório sobre esses textos, como a recuperação pela memória, pela leitura das ilustrações etc. ) Aprendizagens específicas Ler, mesmo que de forma não convencional, textos poéticos, listas, trechos de contos, conhecendo de memória alguns deles.
Aprimorar a comunicação, tornando -a mais objetiva e estruturada, participando de diferentes situações de leitura, com variados gêneros textuais, feita pelos \textbf{adultos} ou por si mesma ( mesmo que de forma não-convencional ).
Ampliar seu repertório cultural literário, desenvolvendo sensibilidade, criatividade, imaginação e gosto pela leitura por meio dos tipos de textos que tem acesso.
Utilizar a estratégia de antecipação do significado dos textos, para confirmação ou não de sua hipótese. \textbf{Sugestão} de experiências - \textbf{Conversar} sobre seus textos preferidos com outras pessoas. - Escolher \textbf{livros} da \textbf{instituição} para ler em casa. - Organizar caixas ou cantos com diversos materiais escritos. - Ler diferentes tipos de textos, mesmo que não convencionalmente. - Ler, mesmo que não convencionalmente, para os \textbf{colegas} diferentes tipos de textos. - Ler, mesmo que não convencionalmente, em situações \textbf{lúdicas} envolvendo o jogo simbólico. 25.2.9.
Objetivo de aprendizagem e desenvolvimento : ( EI03EF09 ) - Levantar hipóteses em relação à linguagem escrita, realizando \textbf{registros} de palavras e textos, por meio de escrita espontânea.
Aprendizagens específicas - Identificar e reconhecer o próprio nome e dos \textbf{colegas} da turma, realizando leitura em situações que se fizer necessário. - Escrever, convencionalmente e não convencionalmente, o seu nome completo, sabendo identificá -lo nas diversas situações do cotidiano. - Manifestar intencionalidade em se comunicar por meio da escrita, expressando suas ideias e pensamentos em situações de uso social. - Conhecer, progressivamente, as diversas funções sociais da linguagem escrita, utilizando variados tipos de textos e gêneros textuais, nas situações cotidianas, em que o uso destes se faça necessário. - Estabelecer relação entre fala e escrita, compreendendo como esta se organiza, fazendo o uso desse conhecimento nas diversas práticas de leitura e escrita do cotidiano. \textbf{Sugestões} experiências - Grafar o próprio nome, para identificar trabalhos, materiais pessoais e jogos. - Elaborar avisos, pedidos e recados. - Ler diferentes tipos de textos. - Produzir diferentes tipos de textos. - Identificar as escritas que existem no espaço escolar. 25.3.
Avaliação - \textbf{acompanhamento} do processo de aprendizagem e desenvolvimento das \textbf{crianças} \textbf{pequenas} : Observação \textbf{criteriosa}, \textbf{registro} e análise quanto : - Às diferentes maneiras de expressar ideias e sentimentos. - Ao uso da linguagem oral para se comunicar. - À demonstração de interesse em participar das propostas de leitura e escrita. - À escrita do seu nome e dos \textbf{colegas}. - Ao interesse por escrever seu próprio nome e as pistas que utiliza para escrevê -lo. - Às estratégias e procedimentos que utiliza para leitura e produção de textos. - Ao interesse por escrever e compartilhar suas produções escritas. - À familiaridade com os diversos gêneros literários, autores, características e suportes. - Ao uso de \textbf{livros}, materiais impressos, sala de leitura e \textbf{biblioteca}. - À familiaridade com materiais e as novas tecnologias para produção da escrita e seus diferentes usos. - À participação na produção de textos orais em diferentes situações. - À ampliação do tempo de \textbf{escuta}, \textbf{interesse} pelo que \textbf{ouve} e se faz perguntas e comentários sobre o assunto. - Ao interesse pela apreciação de textos poéticos e narrativos. - Ao interesse por recitar textos de memória, utilizando -os nos momentos de jogos e \textbf{brincadeiras}. - Ao interesse para ler diferentes textos, ainda que de forma não convencional. - À interpretação feita sobre o assunto trabalhado. - Ao interesse em \textbf{ouvir} a leitura de diversos textos. - À atitude de participação oral durante as atividades.

% matched lemmas: acompanhamento, adulto, ajuda, apoiar, atento, bebê, biblioteca, brincadeira, brincar, brinquedo, colega, contar, conversar, criança, criterioso, cuidado, curiosidade, demonstrar, desejo, dramatizar, dramático, escuta, escutar, estímulo, foto, gesto, hábito, imersão, indicativo, infantil, instituição, interessar, inventar, livro, lúdico, manusear, manuseio, mês, nascimento, ouvir, pequeno, percepção, personagem, plástico, registro, roda, som, sugestão
\end{textsample}
